\begin{table*}[t]
\centering\footnotesize
\setlength{\tabcolsep}{4pt}
\begin{tabular}{@{}c*{9}{r}@{}}
\toprule
& \multicolumn{3}{c}{$Q\cdot P=0$} & \multicolumn{3}{c}{$Q\cdot P=1$} & \multicolumn{3}{c}{$Q\cdot P=2$}\\
\cmidrule(lr){2-4}\cmidrule(lr){5-7}\cmidrule(lr){8-10}
$Q^2$ & $P^2=0$ & $P^2=2$ & $P^2=4$ & $P^2=0$ & $P^2=2$ & $P^2=4$ & $P^2=0$ & $P^2=2$ & $P^2=4$ \\
\midrule
0 & 2 & 64 & 816 & 0 & 0 & -8 & 0 & 0 & 0 \\
1 & -12 & -224 & -1\,248 & 1 & 96 & 1\,968 & 0 & 0 & -12 \\
2 & 48 & 1\,152 & 18\,240 & -8 & -256 & 840 & 0 & 64 & 2\,592 \\
3 & -168 & -3\,392 & -10\,320 & 37 & 1\,376 & 34\,656 & -2 & -224 & 2\,432 \\
4 & 528 & 11\,520 & 200\,736 & -136 & -3\,840 & 16\,632 & 16 & 1\,152 & 43\,392 \\
5 & -1\,512 & -30\,336 & -55\,424 & 439 & 13\,152 & 343\,152 & -72 & -3\,392 & 33\,720 \\
6 & 4\,032 & 83\,968 & 1\,544\,832 & -1\,288 & -33\,536 & 171\,152 & 256 & 11\,520 & 414\,336 \\
7 & -10\,176 & -202\,560 & -179\,022 & 3\,503 & 92\,928 & 2\,476\,752 & -806 & -30\,336 & 302\,400 \\
8 & 24\,528 & 496\,512 & 9\,480\,000 & -8\,968 & -220\,672 & 1\,265\,256 & 2\,320 & 83\,968 & 2\,926\,080 \\
9 & -56\,796 & -1\,118\,496 & -155\,232 & 21\,854 & 540\,416 & 14\,545\,584 & -6\,200 & -202\,560 & 2\,049\,968 \\
10 & 127\,008 & 2\,521\,600 & 49\,523\,328 & -51\,080 & -1\,204\,992 & 7\,558\,560 & 15\,616 & 496\,512 & 16\,919\,712 \\
11 & -275\,544 & -5\,374\,656 & 2\,684\,560 & 115\,154 & 2\,711\,616 & 73\,540\,080 & -37\,508 & -1\,118\,496 & 11\,568\,000 \\
12 & 581\,952 & 11\,389\,440 & 228\,872\,064 & -251\,528 & -5\,741\,824 & 38\,736\,600 & 86\,544 & 2\,521\,600 & 84\,554\,880 \\
\bottomrule
\end{tabular}
\caption{The index $d(Q,P)$ of quarter-BPS dyons in the $\ZZ_2$ toroidal orbifold of type IIB on $T^6$, from the product formula for $1/\Phit_2$ with the twisted elliptic genera \eqref{eq:tor2genera}, in the chamber $\operatorname{Im} v<0$ and with the normalisation of Remark~\ref{rem:normalisation}. The entries agree with Tables~17, 18 and~19 of \cite{CD19} in their common range $Q^2\le 12$, except the entry $Q^2=4$, $P^2=2$, $Q\cdot P=0$, printed there as $-11\,520$: the product formula gives $+11\,520$, and the sign rule $(-1)^{Q^2}$ for $P^2=2$ stated in the same section of \cite{CD19} requires the positive value. The column $P^2=0$ lies outside the single-centred cone; its sign is $(-1)^{Q^2}$. The negative entries inside the cone are those of Table~\ref{tab:violations}.}
\label{tab:tor2}
\end{table*}
